% ===========================================================================
%  Chapitre 14 — Transformations du plan et isométries
%  PE 2026, composante GÉOMÉTRIE
% ===========================================================================
\bmtchapitre{Transformations du plan et isométries}
  {Construire l'image d'une figure par une transformation plane, composer deux
   transformations, catégoriser une isométrie et la décomposer en produit de
   réflexions.}
  {Géométrie \textbullet\ environ 10 heures}

Jusqu'ici, vous avez utilisé les transformations une par une : une translation
pour déplacer une figure, une rotation pour la faire tourner, une symétrie
pour la retourner. Ce chapitre change de point de vue et les considère comme
des objets qu'on peut composer, comme on compose deux fonctions. La question
devient alors : qu'obtient-on en enchaînant deux transformations ?

La réponse est remarquable. Quatre transformations seulement conservent les
distances — translation, rotation, réflexion, symétrie glissée — et toute
composée retombe sur l'une des quatre. Mieux : chacune s'écrit comme un
produit de réflexions. Reconnaître une isométrie, c'est donc identifier deux
axes, et c'est exactement ce que demande le problème de géométrie du
baccalauréat.

% ---------------------------------------------------------------------------
\begin{revision}
Géométrie de première et chapitre précédent.

\begin{enumerate}[label=\textbf{\arabic*.}, leftmargin=*, itemsep=0.35em]
  \item Construire l'image d'un triangle par la translation de vecteur
        $\vecu$, puis par la symétrie de centre $O$.
  \item Rappeler la définition d'une rotation de centre $O$ et d'angle
        $\theta$.
  \item Qu'est-ce que la médiatrice d'un segment ? Donner sa caractérisation
        par les distances.
  \item Rappeler la définition d'une homothétie de centre $O$ et de rapport
        $k$.
  \item Que conserve une homothétie de rapport $2$ : les longueurs, les
        angles, les aires ?
\end{enumerate}
\end{revision}

\begin{automatismes}
Sans calculatrice, sans brouillon.

\begin{enumerate}[label=\textbf{\alph*)}, leftmargin=*, itemsep=0.25em]
  \item Image de $M$ par $\trans{u}$ si $\vect{MM'} = \vecu$
  \item Composée de deux translations de vecteurs $\vecu$ et $\vecv$
  \item Angle de la composée de deux rotations d'angles
        $\frac{\pi}{3}$ et $\frac{\pi}{6}$
  \item Composée d'une réflexion par elle-même
  \item Rapport de la composée de deux homothéties de rapports $2$ et $3$
  \item Une rotation d'angle $\pi$ est aussi\dots
  \item Points invariants d'une réflexion d'axe $(D)$
  \item Points invariants d'une translation de vecteur non nul
\end{enumerate}
\end{automatismes}

\begin{corrige}[automatismes]
a) le point $M'$ \quad b) la translation de vecteur $\vecu+\vecv$ \quad
c) $\frac{\pi}{2}$ \quad d) l'identité \quad e) $6$ \quad
f) une symétrie centrale \quad g) tous les points de $(D)$ \quad
h) aucun.
\end{corrige}

% ===========================================================================
\section{Les transformations usuelles du plan}

\begin{definition}[transformation]
Une \textbf{transformation} du plan est une application bijective du plan dans
lui-même. Un point $M$ tel que $f(M) = M$ est dit \textbf{invariant}.
\end{definition}

\begin{definition}[les six transformations du programme]
\begin{itemize}[leftmargin=*, itemsep=0.15em]
  \item \textbf{Translation} $\trans{u}$ : $\vect{MM'} = \vecu$.
  \item \textbf{Homothétie} $\homo{O}{k}$ : $\vect{OM'} = k\,\vect{OM}$, avec
        $k \neq 0$.
  \item \textbf{Rotation} $\rot{O}{\theta}$ : $OM' = OM$ et
        $\left(\vect{OM},\vect{OM'}\right) = \theta$.
  \item \textbf{Réflexion} $s_{(D)}$ : symétrie orthogonale d'axe $(D)$.
  \item \textbf{Symétrie centrale} $s_{O}$ : c'est $\rot{O}{\pi}$, ou encore
        $\homo{O}{-1}$.
  \item \textbf{Symétrie glissée} : composée
        $\trans{u} \circ s_{(D)}$ où $\vecu$ est un vecteur directeur non nul
        de $(D)$.
\end{itemize}
\end{definition}

\begin{propriete}[ce que chacune conserve]
Toutes conservent l'alignement, le parallélisme, les milieux et les rapports
de longueurs. Les distances sont conservées par la translation, la rotation,
la réflexion et la symétrie glissée ; l'homothétie de rapport $k$ les
multiplie par $\abs k$. L'orientation des angles est conservée par la
translation, la rotation et l'homothétie, et inversée par la réflexion et la
symétrie glissée.
\end{propriete}

\begin{avertissement}
Une symétrie glissée n'a \emph{aucun} point invariant, alors qu'une réflexion
en a toute une droite. C'est le critère qui les distingue immédiatement, et
c'est la première chose à vérifier quand on identifie un antidéplacement.
\end{avertissement}

% ===========================================================================
\section{Composition de deux transformations}

\begin{propriete}[quelques composées immédiates]
\begin{enumerate}[label=\textbf{\arabic*.}, leftmargin=*, itemsep=0.15em]
  \item $\trans{u} \circ \trans{v} = \trans{u+v}$.
  \item $\rot{O}{\theta} \circ \rot{O}{\theta'} = \rot{O}{\theta+\theta'}$.
  \item $\homo{O}{k} \circ \homo{O}{k'} = \homo{O}{kk'}$.
  \item La composée de deux rotations de centres distincts et d'angles
        $\theta$, $\theta'$ est une rotation d'angle $\theta+\theta'$ si
        $\theta+\theta' \neq 0 \ [2\pi]$, et une translation sinon.
\end{enumerate}
\end{propriete}

\begin{avertissement}
La composition n'est pas commutative. Les transformations
$\trans{u} \circ \rot{O}{\theta}$ et $\rot{O}{\theta} \circ \trans{u}$ sont
toutes deux des rotations d'angle $\theta$, mais leurs centres diffèrent. Le
problème de baccalauréat exploite régulièrement cette distinction en
demandant successivement $f = t \circ r$ et $g = r \circ t$.
\end{avertissement}

% ===========================================================================
\section{Isométries : déplacements et antidéplacements}

\begin{definition}[isométrie]
Une \textbf{isométrie} est une transformation du plan qui conserve les
distances : $f(M)f(N) = MN$ pour tous points $M$ et $N$. Elle est appelée
\textbf{déplacement} si elle conserve les angles orientés,
\textbf{antidéplacement} si elle les change en leurs opposés.
\end{definition}

\begin{theoreme}[classification des isométries du plan]
Les déplacements du plan sont exactement les translations et les rotations.
Les antidéplacements sont exactement les réflexions et les symétries
glissées.
\end{theoreme}

\begin{propriete}[reconnaissance par les points invariants]
\begin{center}
\setlength{\tabcolsep}{9pt}
\renewcommand{\arraystretch}{1.25}
\begin{tabular}{@{}lll@{}}
\toprule
\textsf{\footnotesize\bfseries POINTS INVARIANTS} &
\textsf{\footnotesize\bfseries DÉPLACEMENT} &
\textsf{\footnotesize\bfseries ANTIDÉPLACEMENT} \\
\midrule
tout le plan & identité & — \\
une droite   & — & réflexion \\
un seul point & rotation & — \\
aucun & translation & symétrie glissée \\
\bottomrule
\end{tabular}
\end{center}
\end{propriete}

\begin{methode}{identifier une isométrie donnée par des images}
Trois questions, dans cet ordre.
\begin{enumerate}[label=\textbf{\arabic*.}, leftmargin=*, itemsep=0.15em]
  \item \emph{Déplacement ou antidéplacement ?} On compare l'orientation d'un
        triangle et celle de son image.
  \item \emph{Y a-t-il un point invariant ?} On résout $f(M) = M$, ou l'on
        raisonne sur la figure.
  \item \emph{Éléments caractéristiques.} Vecteur pour une translation ;
        centre et angle pour une rotation ; axe pour une réflexion ; axe et
        vecteur pour une symétrie glissée.
\end{enumerate}
\end{methode}

% ===========================================================================
\section{Composée de deux réflexions}

\begin{theoreme}[axes parallèles]
Si $(D)$ et $(D')$ sont parallèles, alors
\[
  s_{(D')} \circ s_{(D)} = \trans{2\,\vect{HH'}},
\]
où $H \in (D)$ et $H'$ est son projeté orthogonal sur $(D')$. C'est une
translation de vecteur orthogonal aux deux axes, de norme le double de leur
distance.
\end{theoreme}

\begin{theoreme}[axes sécants]
Si $(D)$ et $(D')$ se coupent en $O$ en formant l'angle
$\theta = \left((D),(D')\right)$, alors
\[
  s_{(D')} \circ s_{(D)} = \rot{O}{2\theta}.
\]
\end{theoreme}

\begin{visualisation}[deux miroirs, une rotation]
\begin{center}
\begin{tikzpicture}[x=1.1cm, y=1.1cm, font=\footnotesize]
  \coordinate (O) at (0,0);
  \draw[bmtGris, line width=1pt] (-2.6,0) -- (2.6,0) node[anchor=west] {$(D)$};
  \draw[bmtGris, line width=1pt, rotate=35] (-2.6,0) -- (2.6,0)
    node[anchor=south west] {$(D')$};
  \filldraw[bmtEncre] (O) circle (1.8pt) node[anchor=north east] {$O$};
  \coordinate (M) at (55:2.2);
  \coordinate (M1) at (-55:2.2);
  \coordinate (M2) at (125:2.2);
  \filldraw[bmtIndigo] (M) circle (2pt) node[anchor=south west] {$M$};
  \filldraw[bmtVetiver] (M1) circle (2pt) node[anchor=north west] {$M_{1}$};
  \filldraw[bmtLaterite] (M2) circle (2pt) node[anchor=south east] {$M'$};
  \draw[bmtVetiver, dashed] (M) -- (M1);
  \draw[bmtLaterite, dashed] (M1) -- (M2);
  \draw[bmtLaterite, line width=0.9pt, -{Stealth[length=2.4mm]}]
    (55:1.1) arc (55:125:1.1);
  \node[bmtLaterite, anchor=south] at (90:1.25) {$2\theta$};
  \draw[bmtGris, line width=0.8pt] (0:0.62) arc (0:35:0.62);
  \node[bmtGris, anchor=west] at (17:0.75) {$\theta$};
\end{tikzpicture}
\end{center}

On réfléchit $M$ dans le premier miroir, puis l'image dans le second. Chaque
réflexion inverse l'orientation ; deux réflexions la rétablissent, et le
résultat est donc un déplacement. Le point $O$, commun aux deux axes, est
invariant : c'est une rotation. L'angle balayé vaut le double de l'angle des
axes — la figure le montre sans calcul.
\end{visualisation}

\begin{demonstration}
Le point $O$ appartient aux deux axes, donc il est invariant par chacune des
réflexions, donc par la composée. Celle-ci conserve les distances et
l'orientation — deux inversions se compensent — : c'est un déplacement fixant
un point, donc une rotation de centre $O$. Pour l'angle, on suit un point $M$
de $(D)$ : son image par $s_{(D)}$ est $M$ lui-même, et
$\left(\vect{OM},\vect{OM'}\right) = 2\theta$ par définition de la réflexion
d'axe $(D')$.
\end{demonstration}

% ===========================================================================
\section{Décomposition en produit de réflexions}

\begin{theoreme}[décomposition d'une translation]
Toute translation $\trans{u}$ se décompose en produit de deux réflexions
d'axes parallèles, perpendiculaires à $\vecu$ et distants de
$\frac{\norme{\vecu}}{2}$. \emph{L'un des deux axes peut être choisi
librement} parmi les droites perpendiculaires à $\vecu$ ; l'autre est alors
imposé.
\end{theoreme}

\begin{theoreme}[décomposition d'une rotation]
Toute rotation $\rot{O}{\theta}$ se décompose en produit de deux réflexions
d'axes sécants en $O$ formant un angle $\frac{\theta}{2}$. \emph{L'un des deux
axes peut être choisi librement} parmi les droites passant par $O$.
\end{theoreme}

\begin{methode}{composer deux isométries par décomposition}
C'est la technique reine du chapitre, et elle repose sur la liberté de choix
de l'un des axes.
\begin{enumerate}[label=\textbf{\arabic*.}, leftmargin=*, itemsep=0.15em]
  \item Écrire chacune des deux transformations comme produit de deux
        réflexions.
  \item Choisir l'axe libre de chacune de sorte que les deux réflexions
        \emph{intérieures} soient identiques.
  \item Elles se simplifient — une réflexion composée avec elle-même est
        l'identité —, et il ne reste qu'un produit de deux réflexions.
  \item Conclure : axes parallèles donnent une translation, axes sécants une
        rotation.
\end{enumerate}
\end{methode}

\begin{exemple}[la composée d'une rotation et d'une translation]
Soit $r$ la rotation de centre $C$ et d'angle $\frac{\pi}{2}$, et $t$ la
translation de vecteur $\vect{CA}$. Cherchons la nature de
$f = t \circ r$.

La composée conserve les distances et l'orientation : c'est un déplacement.
Son angle est celui de $r$, soit $\frac{\pi}{2}$, non nul : c'est donc une
rotation d'angle $\frac{\pi}{2}$, et il reste à trouver son centre.

Décomposons. Soit $(\Delta)$ la droite passant par $C$ et perpendiculaire à
$\vect{CA}$. On écrit
\[
  r = s_{(\Delta)} \circ s_{(\Delta_{1})},
  \qquad
  t = s_{(\Delta_{2})} \circ s_{(\Delta)},
\]
où $(\Delta_{1})$ passe par $C$ en formant avec $(\Delta)$ l'angle
$-\frac{\pi}{4}$, et où $(\Delta_{2})$ est parallèle à $(\Delta)$, à la
distance $\frac{CA}{2}$. Alors
\[
  f = s_{(\Delta_{2})} \circ \underbrace{s_{(\Delta)} \circ s_{(\Delta)}}_{
  \text{identité}} \circ\, s_{(\Delta_{1})}
  = s_{(\Delta_{2})} \circ s_{(\Delta_{1})},
\]
produit de deux réflexions d'axes sécants : le centre de $f$ est le point
d'intersection de $(\Delta_{1})$ et $(\Delta_{2})$, et l'angle vaut le double
de l'angle des axes, soit $\frac{\pi}{2}$.
\end{exemple}

\begin{entrainement}[composer et décomposer]
\exo[\niveau{1}] Déterminer la nature et les éléments de
$\rot{O}{\pi/3} \circ \rot{O}{2\pi/3}$.

\exo[\niveau{2}] $(D)$ et $(D')$ sont parallèles et distantes de $3$~cm.
Déterminer $s_{(D')} \circ s_{(D)}$.

\exo[\niveau{2}] $(D)$ et $(D')$ se coupent en $O$ avec un angle de
$30^{\circ}$. Déterminer $s_{(D')} \circ s_{(D)}$.

\exo[\niveau{2}] Décomposer la translation de vecteur $\vecu$ en deux
réflexions dont l'un des axes est une droite donnée perpendiculaire à
$\vecu$.

\exo[\niveau{3}] Soient $r_{1}$ de centre $A$ et d'angle $\frac{\pi}{2}$, et
$r_{2}$ de centre $O$ et d'angle $\frac{\pi}{2}$. Montrer que
$r_{1} \circ r_{2}$ est une symétrie centrale et en déterminer le centre par
décomposition.
\end{entrainement}

\begin{corrige}
\textbf{1.} Deux rotations de même centre se composent en additionnant les
angles : $\frac{\pi}{3}+\frac{2\pi}{3} = \pi$. C'est donc la rotation de
centre $O$ et d'angle $\pi$, c'est-à-dire la symétrie centrale de centre $O$.

\textbf{2.} La composée de deux réflexions d'axes parallèles est la
translation dont le vecteur est perpendiculaire aux axes, dirigé de $(D)$ vers
$(D')$, et de norme le double de leur distance, soit $6$~cm.

\textbf{3.} La composée de deux réflexions d'axes sécants en $O$ est la
rotation de centre $O$ et d'angle le double de l'angle orienté qui mène du
premier axe au second, soit $60^{\circ}$.

\textbf{4.} Soit $(D)$ la droite donnée, perpendiculaire à $\vecu$. En notant
$(D')$ son image par la translation de vecteur $\frac12\vecu$, les deux axes
sont parallèles et distants de $\frac12\norme{\vecu}$ : la composée
$s_{(D')} \circ s_{(D)}$ est bien la translation de vecteur $\vecu$.

\textbf{5.} Les angles s'ajoutent : $\frac{\pi}{2}+\frac{\pi}{2} = \pi$, donc
la composée est une rotation d'angle $\pi$, c'est-à-dire une symétrie
centrale. Pour trouver son centre, on décompose chaque rotation en deux
réflexions en choisissant la droite $(OA)$ comme axe commun :
$r_{2} = s_{(OA)} \circ s_{a}$ et $r_{1} = s_{d} \circ s_{(OA)}$, où $a$ passe
par $O$ et $d$ par $A$, chacune faisant un angle de $\frac{\pi}{4}$ avec
$(OA)$, mais de part et d'autre. Alors
$r_{1} \circ r_{2} = s_{d} \circ s_{a}$, et le centre cherché est le point
d'intersection de $a$ et de $d$. En repérant $O(0\,;0)$ et $A(a\,;0)$, ce
point a pour coordonnées $\left(\frac a2\,;-\frac a2\right)$ ; on le vérifie
en constatant qu'il est bien le seul point invariant de la composée.
\end{corrige}

% ===========================================================================
\section{Isométries admettant un point invariant}

\begin{theoreme}[isométries fixant un point]
Une isométrie du plan admettant au moins un point invariant est soit
l'identité, soit une rotation de centre ce point, soit une réflexion dont
l'axe passe par ce point.
\end{theoreme}

\begin{demonstration}
Soit $f$ une isométrie fixant $O$. Si $f$ est un déplacement, elle est
translation ou rotation ; une translation fixant un point est l'identité, donc
$f$ est l'identité ou une rotation de centre $O$. Si $f$ est un
antidéplacement, elle est réflexion ou symétrie glissée ; une symétrie glissée
n'a aucun point invariant, donc $f$ est une réflexion, et son axe contient
$O$.
\end{demonstration}

\begin{propriete}[une isométrie est déterminée par trois points]
Une isométrie du plan est entièrement déterminée par les images de trois
points non alignés. En particulier, deux isométries qui coïncident sur trois
points non alignés sont égales.
\end{propriete}

% ===========================================================================
\begin{annales}
Les transformations occupent la première partie du problème de géométrie, et
la décomposition en réflexions y est demandée presque chaque année. Les
énoncés qui suivent sont repris de sujets réellement posés ; leur provenance
est indiquée à chaque fois.

\exoann{Baccalauréat, Madagascar, série S, session 2021 — problème 1, partie A}
$ABC$ est rectangle en $A$ avec $AB = AC = 3$~cm et
$\left(\vect{AB},\vect{AC}\right) = \frac{\pi}{2}$. On note $r$ la rotation de
centre $C$ et d'angle $\frac{\pi}{2}$, $t$ la translation de vecteur
$\vect{CA}$, et l'on pose $f = t \circ r$. Soient $E$ le symétrique de $B$ par
rapport à $A$, $D$ le barycentre de
$\left\{(A\,;1),(B\,;-1),(C\,;1)\right\}$ et $K$ le milieu de $[ED]$.
\begin{enumerate}[label=\textbf{\alph*)}, leftmargin=*, itemsep=0.15em]
  \item Donner la nature de $f$.
  \item En décomposant $t$ et $r$ en deux symétries orthogonales
        convenablement choisies, déterminer les éléments caractéristiques de
        $f$.
\end{enumerate}

\exoann{Baccalauréat, Madagascar, série S, session 2021 — problème 1, partie A, suite}
Avec les mêmes données, $I$ est le milieu de $[BC]$ et $h$ l'homothétie de
centre $I$ et de rapport $2$. On pose $g = h \circ f$. En déduire la nature et
les éléments caractéristiques de $g$.

\exoann{Baccalauréat, Madagascar, série S, session 2022 — problème 1, partie A}
$ABCD$ est un carré direct de côté $4$~cm. On note $t$ la translation de
vecteur $\vect{AC}$, $t_{1}$ celle de vecteur $\vect{DC}$, $t_{2}$ celle de
vecteur $\vect{AD}$, $s_{(AD)}$ la réflexion d'axe $(AD)$, et
$f = t \circ s_{(AD)}$.
\begin{enumerate}[label=\textbf{\alph*)}, leftmargin=*, itemsep=0.15em]
  \item Justifier que $t = t_{1} \circ t_{2}$.
  \item Décomposer $t_{1}$ en deux symétries orthogonales dont l'un des axes
        est la droite $(AD)$.
  \item En déduire la nature et les éléments caractéristiques de $f$.
\end{enumerate}

\exoann{Baccalauréat, Madagascar, série S, session 2023 — problème 1, partie I}
$ABC$ est rectangle et isocèle en $A$, $AB = 4$~cm. On note $O$ le milieu de
$[BC]$, $r_{1}$ la rotation de centre $A$ et d'angle $\frac{\pi}{2}$,
$r_{2}$ celle de centre $O$ et d'angle $\frac{\pi}{2}$, et $f = r_{1} \circ
r_{2}$.
\begin{enumerate}[label=\textbf{\alph*)}, leftmargin=*, itemsep=0.15em]
  \item Justifier que $f$ est une symétrie centrale.
  \item En décomposant $r_{1}$ et $r_{2}$ en deux symétries orthogonales
        d'axes convenablement choisis, déterminer le centre de $f$.
\end{enumerate}

\exoann{Baccalauréat, Madagascar, série S, session 2026 — problème 1, partie I}
$ABC$ est direct, rectangle en $A$, avec $AC = 4$~cm et $BC = 8$~cm. On note
$R_{C}$ la rotation de centre $C$ et d'angle $\frac{\pi}{3}$, $R_{A}$ celle de
centre $A$ et d'angle $-\frac{\pi}{3}$, et l'on pose $g = R_{C} \circ R_{A}$.
Identifier $g$ en décomposant $R_{A}$ et $R_{C}$ en produits de symétries
orthogonales, puis préciser ses éléments caractéristiques.

\exoann{Baccalauréat blanc, lycée Philibert Tsiranana, Terminale S, 2021-2022 — problème 1, partie B}
Dans le plan orienté, $OAB$ est tel que $OA = OB = a$ et
$\left(\vect{OA},\vect{OB}\right) = \frac{\pi}{2}$. On note $R$ la rotation de
centre $O$ et d'angle $\frac{\pi}{2}$, $T$ la translation de vecteur
$\vect{AB}$, et l'on pose $f = R \circ T$ et $g = T \circ R$.
\begin{enumerate}[label=\textbf{\alph*)}, leftmargin=*, itemsep=0.15em]
  \item Préciser la nature de $f$ et de $g$.
  \item Décomposer $R$ et $T$ en deux réflexions d'axes convenables, puis en
        déduire les éléments caractéristiques de $f$ et de $g$.
  \item Déterminer la nature de $g \circ f$, puis l'image de $A$ par cette
        transformation.
\end{enumerate}

\exoann{Baccalauréat blanc, lycée Philibert Tsiranana, Terminale S, 2022-2023 — problème 1, partie A}
$OAC$ est un triangle équilatéral tel que
$\left(\vect{OA},\vect{OC}\right) = \frac{\pi}{3}$ et $OA = a$. On note
$R_{O}$ la rotation de centre $O$ et d'angle $\frac{\pi}{3}$, $R_{A}$ celle de
centre $A$ et d'angle $\frac{2\pi}{3}$, et $R = R_{O} \circ R_{A}$.
\begin{enumerate}[label=\textbf{\alph*)}, leftmargin=*, itemsep=0.15em]
  \item Montrer que $R$ est une symétrie centrale.
  \item Déterminer $R(A)$ et en déduire le centre de $R$.
  \item Déterminer les droites $(D_{1})$ et $(D_{2})$ telles que
        $R_{O} = s_{(D_{2})} \circ s_{(OA)}$ et
        $R_{A} = s_{(OA)} \circ s_{(D_{1})}$, puis retrouver les résultats
        précédents.
\end{enumerate}

\exoann{Baccalauréat, Mauritanie, série C, session 2022 — exercice 3}
$ABC$ est un triangle équilatéral direct de centre $O$, $D$ est le symétrique
de $O$ par rapport à $(AC)$, et $I$, $J$, $K$ sont les milieux respectifs de
$[AB]$, $[BC]$, $[CA]$.
\begin{enumerate}[label=\textbf{\alph*)}, leftmargin=*, itemsep=0.15em]
  \item Montrer qu'il existe un unique déplacement $f$ transformant $B$ en $C$
        et $O$ en $D$.
  \item Justifier que $f$ est une rotation et préciser ses éléments.
  \item Déterminer la nature et les éléments de
        $\trans{AB} \circ f$.
\end{enumerate}

\exoann{Baccalauréat, Mauritanie, série C, session 2022 — exercice 3, suite}
Avec les mêmes données :
\begin{enumerate}[label=\textbf{\alph*)}, leftmargin=*, itemsep=0.15em]
  \item Montrer qu'il existe un unique antidéplacement $g$ transformant $A$ en
        $C$ et $K$ en $J$.
  \item Prouver que $g$ est une symétrie glissée et donner sa forme réduite.
\end{enumerate}

\exoann{Baccalauréat, Congo-Brazzaville, série C, session 2022 — exercice 2}
$ABC$ est un triangle équilatéral direct et $D$ est le point tel que $ACDB$
soit un losange. On pose $g = s_{(BC)} \circ \trans{CA}$.
\begin{enumerate}[label=\textbf{\alph*)}, leftmargin=*, itemsep=0.15em]
  \item Montrer que $g$ est une symétrie glissée.
  \item Déterminer son vecteur et son axe.
\end{enumerate}
\end{annales}

\begin{corrige}[annales]
\textbf{1. a)} $f$ est composée d'un déplacement et d'une translation : c'est
un déplacement, d'angle celui de $r$, soit $\frac{\pi}{2} \neq 0$. C'est donc
une rotation d'angle $\frac{\pi}{2}$.
\textbf{b)} Soit $(\Delta)$ la perpendiculaire à $\vect{CA}$ passant par $C$.
On écrit $r = s_{(\Delta)} \circ s_{(\Delta_{1})}$ avec $(\Delta_{1})$ passant
par $C$ et formant l'angle $-\frac{\pi}{4}$ avec $(\Delta)$, puis
$t = s_{(\Delta_{2})} \circ s_{(\Delta)}$ avec $(\Delta_{2})$ parallèle à
$(\Delta)$ à la distance $\frac{CA}{2}$. Les réflexions intérieures se
simplifient et $f = s_{(\Delta_{2})} \circ s_{(\Delta_{1})}$ : le centre est
le point d'intersection des deux axes, et l'on vérifie qu'il s'agit de $K$,
milieu de $[ED]$.

\textbf{2.} $g = h \circ f$ compose une rotation d'angle $\frac{\pi}{2}$ et
une homothétie de rapport $2$. La composée multiplie donc les longueurs par
$2$ et fait tourner les directions de $\frac{\pi}{2}$ ; elle possède un unique
point invariant, obtenu en résolvant $g(M) = M$, ou plus rapidement à l'aide
de l'expression complexe de la partie B du sujet. Le chapitre 15 donnera un
nom à cette transformation.

\textbf{3. a)} $\vect{AC} = \vect{AD}+\vect{DC}$, donc
$t = t_{1} \circ t_{2}$, les translations commutant.
\textbf{b)} $t_{1}$, de vecteur $\vect{DC}$ perpendiculaire à $(AD)$, se
décompose en $s_{(\Delta)} \circ s_{(AD)}$ où $(\Delta)$ est parallèle à
$(AD)$, à la distance $\frac{DC}{2} = 2$.
\textbf{c)} Alors
$f = t_{1} \circ t_{2} \circ s_{(AD)}
= s_{(\Delta)} \circ s_{(AD)} \circ t_{2} \circ s_{(AD)}$. Comme $t_{2}$ a
pour vecteur $\vect{AD}$, dirigé par l'axe, la composée
$t_{2} \circ s_{(AD)}$ est une symétrie glissée d'axe $(AD)$ et de vecteur
$\vect{AD}$. Au total, $f$ est une symétrie glissée d'axe $(\Delta)$ et de
vecteur $\vect{AD}$.

\textbf{4. a)} La somme des angles vaut $\frac{\pi}{2}+\frac{\pi}{2} = \pi$ :
$f$ est une rotation d'angle $\pi$, c'est-à-dire une symétrie centrale.
\textbf{b)} On décompose $r_{2} = s_{(AO)} \circ s_{(\Delta_{2})}$ et
$r_{1} = s_{(\Delta_{1})} \circ s_{(AO)}$, où $(AO)$ est l'axe commun choisi.
Après simplification, $f = s_{(\Delta_{1})} \circ s_{(\Delta_{2})}$ : les deux
axes étant perpendiculaires, la composée est la symétrie centrale de centre
leur point d'intersection.

\textbf{5.} La somme des angles vaut $\frac{\pi}{3}-\frac{\pi}{3} = 0$ : $g$
est une translation. En décomposant $R_{A} = s_{(AC)} \circ s_{(\Delta_{1})}$
et $R_{C} = s_{(\Delta_{2})} \circ s_{(AC)}$ le long de l'axe commun $(AC)$,
il vient $g = s_{(\Delta_{2})} \circ s_{(\Delta_{1})}$ avec
$(\Delta_{1})$ et $(\Delta_{2})$ parallèles : le vecteur de la translation est
orthogonal à ces axes et de norme le double de leur distance. On l'obtient
concrètement en calculant $g(A)$.

\textbf{6. a)} $f$ et $g$ sont des déplacements d'angle $\frac{\pi}{2}$ : ce
sont des rotations d'angle $\frac{\pi}{2}$, de centres distincts.
\textbf{b)} On décompose $R$ et $T$ le long d'un axe commun bien choisi : la
perpendiculaire à $\vect{AB}$ passant par $O$ convient. Les centres
s'obtiennent alors comme intersections des axes restants.
\textbf{c)} $g \circ f$ est un déplacement d'angle $\pi$ : c'est une symétrie
centrale, dont le centre se lit sur l'image de $A$.

\textbf{7. a)} La somme des angles vaut $\frac{\pi}{3}+\frac{2\pi}{3} = \pi$ :
$R$ est une symétrie centrale.
\textbf{b)} $R(A) = R_{O}\left(R_{A}(A)\right) = R_{O}(A)$, image de $A$ par
la rotation de centre $O$ et d'angle $\frac{\pi}{3}$ : c'est $C$. Le centre de
la symétrie centrale est donc le milieu de $[AC]$.
\textbf{c)} $(D_{2})$ est l'image de $(OA)$ par la rotation de centre $O$ et
d'angle $\frac{\pi}{6}$ ; $(D_{1})$ est la droite passant par $A$ formant
l'angle $-\frac{\pi}{3}$ avec $(OA)$. La composée redonne bien la symétrie
centrale de centre le milieu de $[AC]$.

\textbf{8. a)} $BC = CD$ car $D$ est le symétrique de $O$ par rapport à
$(AC)$ dans une figure équilatérale : les distances sont conservées, donc un
déplacement envoyant $B$ sur $C$ et $O$ sur $D$ existe, et il est unique
puisqu'un déplacement est déterminé par les images de deux points distincts.
\textbf{b)} Ce déplacement n'est pas une translation, car $\vect{BC}$ et
$\vect{OD}$ ne sont pas égaux : c'est une rotation. Son angle est
$\left(\vect{BO},\vect{CD}\right)$ et son centre le point d'intersection des
médiatrices de $[BC]$ et de $[OD]$.
\textbf{c)} $\trans{AB} \circ f$ est un déplacement de même angle que $f$ :
c'est une rotation de même angle, de centre déplacé.

\textbf{9. a)} Un antidéplacement est déterminé par les images de deux points
distincts, d'où l'existence et l'unicité.
\textbf{b)} Si $g$ avait un point invariant, ce serait une réflexion, et son
axe serait à la fois la médiatrice de $[AC]$ et celle de $[KJ]$ — ce qui n'est
pas le cas. Donc $g$ est une symétrie glissée. Sa forme réduite s'obtient en
prenant pour axe la droite joignant les milieux de $[AC]$ et de $[KJ]$, et
pour vecteur le double du vecteur joignant $A$ au projeté de $A$ sur cet axe.

\textbf{10. a)} $g$ est composée d'une réflexion et d'une translation :
c'est un antidéplacement. Il n'a pas de point invariant, car le vecteur
$\vect{CA}$ n'est pas orthogonal à $(BC)$ : c'est donc une symétrie glissée.
\textbf{b)} On décompose $\vect{CA}$ en une composante $\vecu$ parallèle à
$(BC)$ et une composante $\vecv$ orthogonale. La composante orthogonale
déplace l'axe de $\frac{\norme{\vecv}}{2}$, la composante parallèle donne le
vecteur de glissement : l'axe est la parallèle à $(BC)$ passant par le milieu
de la projection, et le vecteur vaut $\vecu$.
\end{corrige}

% ===========================================================================
\begin{jecherche}[le motif d'un lamba]
Un tisserand compose un motif sur une bande de lamba. Le motif de base est un
triangle $T$, et la bande est obtenue en répétant deux opérations : la
réflexion $s$ d'axe la droite médiane horizontale de la bande, et la
translation $t$ de vecteur horizontal $\vecu$ de longueur $6$~cm.

\begin{enumerate}[label=\textbf{\arabic*.}, leftmargin=*, itemsep=0.3em]
  \item Quelle est la nature de la transformation $t \circ s$ ? Justifier.
  \item Déterminer la nature de $(t \circ s) \circ (t \circ s)$.
  \item Le tisserand veut que le motif se reproduise à l'identique tous les
        $12$~cm. Est-ce compatible avec la question précédente ?
  \item Quelle est la nature de $s \circ t$ ? Cette transformation est-elle
        égale à $t \circ s$ ?
  \item On ajoute une seconde réflexion $s'$, d'axe vertical. Déterminer la
        nature de $s' \circ s$.
  \item Le tisserand souhaite obtenir une rotation d'un quart de tour. Quels
        axes doit-il choisir ?
\end{enumerate}
\end{jecherche}

\begin{corrige}[le motif d'un lamba]
\textbf{1.} Le vecteur $\vecu$ est parallèle à l'axe de $s$ : la composée
$t \circ s$ est donc une symétrie glissée, d'axe la médiane et de vecteur
$\vecu$.

\textbf{2.} $(t \circ s)^{2} = t \circ s \circ t \circ s$. Comme $\vecu$ est
parallèle à l'axe, $s$ et $t$ commutent, et il reste
$t \circ t \circ s \circ s = t^{2}$ : c'est la translation de vecteur
$2\vecu$, de longueur $12$~cm.

\textbf{3.} Oui, et c'est même exactement ce qui est demandé : appliquer deux
fois la symétrie glissée revient à translater de $12$~cm, le motif se
retrouvant dans son orientation d'origine.

\textbf{4.} $s \circ t$ est également une symétrie glissée de même axe et de
même vecteur. Ici les deux composées sont égales, précisément parce que
$\vecu$ dirige l'axe — c'est le seul cas où réflexion et translation
commutent.

\textbf{5.} Deux réflexions d'axes perpendiculaires donnent une rotation
d'angle $2 \times \frac{\pi}{2} = \pi$ : c'est la symétrie centrale de centre
le point d'intersection des deux axes.

\textbf{6.} Une rotation d'un quart de tour est le produit de deux réflexions
dont les axes forment un angle de $\frac{\pi}{4}$, soit $45^{\circ}$. Le
tisserand doit donc croiser ses deux axes à quarante-cinq degrés.
\end{corrige}

% ===========================================================================
\begin{jemeteste}
Répondre par vrai ou faux, en justifiant en une ligne.

\begin{enumerate}[label=\textbf{\arabic*.}, leftmargin=*, itemsep=0.28em]
  \item Une homothétie de rapport $3$ est une isométrie.
  \item La composée de deux réflexions est toujours une rotation.
  \item Une symétrie glissée possède exactement un point invariant.
  \item La composée de deux rotations est toujours une rotation.
  \item Un déplacement sans point invariant est une translation.
  \item Une isométrie est déterminée par les images de deux points distincts.
\end{enumerate}
\end{jemeteste}

\begin{corrige}[je me teste]
\textbf{1.} Faux — elle multiplie les distances par $3$.
\textbf{2.} Faux — c'est une rotation si les axes sont sécants, une
translation s'ils sont parallèles.
\textbf{3.} Faux — elle n'en a aucun.
\textbf{4.} Faux — si la somme des angles est nulle modulo $2\pi$, on obtient
une translation.
\textbf{5.} Vrai — d'après la classification.
\textbf{6.} Faux — deux points ne suffisent pas : il reste deux isométries
possibles, l'une directe et l'autre indirecte. Il en faut trois, non alignés.
\end{corrige}

% ===========================================================================
\begin{commeaubac}
\bmtsource{Baccalauréat, série S, session 2023 — problème 1, partie I}
Dans le plan orienté, $ABC$ est rectangle et isocèle en $A$ tel que
$AB = 4$~cm et $\left(\vect{AB},\vect{AC}\right) = \frac{\pi}{2}$. On désigne
par $O$ le milieu de $[BC]$ et par $E$ celui de $[AC]$. Soient $r_{1}$ la
rotation de centre $A$ et d'angle $\frac{\pi}{2}$, $r_{2}$ celle de centre $O$
et d'angle $\frac{\pi}{2}$, et $f = r_{1} \circ r_{2}$.

\begin{enumerate}[label=\textbf{\arabic*.}, leftmargin=*, itemsep=0.25em]
  \item Justifier que $f$ est une symétrie centrale. \hfill\textup{(0,5 pt)}
  \item En décomposant $r_{1}$ et $r_{2}$ en deux symétries orthogonales
        d'axes convenablement choisis, déterminer le centre de $f$.
        \hfill\textup{(0,75 pt)}
\end{enumerate}
\end{commeaubac}

\begin{corrige}[comme au baccalauréat]
\textbf{1.} $f$ est composée de deux déplacements : c'est un déplacement, dont
l'angle est la somme des angles, soit $\frac{\pi}{2}+\frac{\pi}{2} = \pi$.
Un déplacement d'angle $\pi$ est une rotation d'angle $\pi$, c'est-à-dire une
symétrie centrale.

\textbf{2.} Choisissons la droite $(AO)$ comme axe commun. On écrit
$r_{2} = s_{(AO)} \circ s_{(\Delta_{2})}$, où $(\Delta_{2})$ passe par $O$ et
forme avec $(AO)$ l'angle $-\frac{\pi}{4}$, puis
$r_{1} = s_{(\Delta_{1})} \circ s_{(AO)}$, où $(\Delta_{1})$ passe par $A$ et
forme avec $(AO)$ l'angle $\frac{\pi}{4}$. Alors
\[
  f = s_{(\Delta_{1})} \circ \underbrace{s_{(AO)} \circ s_{(AO)}}_{
  \text{identité}} \circ\, s_{(\Delta_{2})}
  = s_{(\Delta_{1})} \circ s_{(\Delta_{2})} .
\]
Les deux axes forment entre eux un angle de $\frac{\pi}{2}$ : la composée est
bien une symétrie centrale, de centre le point d'intersection de
$(\Delta_{1})$ et $(\Delta_{2})$.
\end{corrige}

% ===========================================================================
\begin{concours}
\bmtsource{D'après un concours d'entrée en école d'ingénieurs}
Dans le plan orienté, on considère trois droites $(D_{1})$, $(D_{2})$ et
$(D_{3})$ concourantes en un point $O$.

\begin{enumerate}[label=\textbf{\arabic*.}, leftmargin=*, itemsep=0.25em]
  \item Montrer que
        $s_{(D_{3})} \circ s_{(D_{2})} \circ s_{(D_{1})}$ est une réflexion,
        et déterminer son axe.
  \item En déduire que le produit d'un nombre impair de réflexions d'axes
        concourants est une réflexion.
  \item Que devient l'énoncé si les trois droites sont parallèles ?
\end{enumerate}
\end{concours}

\begin{corrige}[vers le concours]
\textbf{1.} Le produit $s_{(D_{2})} \circ s_{(D_{1})}$ est la rotation de
centre $O$ et d'angle $2\theta$, où $\theta$ est l'angle des deux axes. Cette
rotation se décompose en $s_{(D_{3})} \circ s_{(\Delta)}$, où $(\Delta)$ passe
par $O$ et est choisie de sorte que l'angle de $(\Delta)$ à $(D_{3})$ vaille
$\theta$. Alors
$s_{(D_{3})} \circ s_{(D_{2})} \circ s_{(D_{1})}
= s_{(D_{3})} \circ s_{(D_{3})} \circ s_{(\Delta)} = s_{(\Delta)}$ : c'est la
réflexion d'axe $(\Delta)$.

\textbf{2.} Une récurrence : deux réflexions consécutives se remplacent par
une rotation de centre $O$, laquelle se recombine avec la suivante en une
réflexion. Le nombre de réflexions diminue de deux à chaque étape, et la
parité impaire est conservée jusqu'à ce qu'il n'en reste qu'une.

\textbf{3.} Si les axes sont parallèles, deux réflexions donnent une
translation de vecteur orthogonal aux axes ; composée avec la troisième
réflexion, elle donne encore une réflexion d'axe parallèle aux trois autres.
La conclusion est la même, le centre commun étant remplacé par une direction
commune.
\end{corrige}

% ===========================================================================
\begin{notepourlaclasse}
La difficulté de ce chapitre n'est pas technique, elle tient à la figure. Un
élève qui trace mal ne trouve rien. Une séance entière consacrée aux
constructions à la règle et au compas, avant toute théorie, est rentable :
image d'un point par une rotation, par une réflexion, par une symétrie
glissée, à la main, dix fois.

La décomposition en réflexions est le cœur du chapitre et le point le moins
bien compris. Insistez sur la \emph{liberté de choix} de l'un des axes :
c'est elle, et elle seule, qui permet de faire disparaître les deux
réflexions intérieures. Tant que cette liberté n'est pas perçue, la méthode
paraît magique et n'est pas réutilisée.

Un dernier conseil : faites toujours annoncer la nature de la transformation
avant d'en chercher les éléments. Somme des angles, conservation ou inversion
de l'orientation, existence d'un point invariant — trois questions, et la
nature est fixée. Chercher un centre avant de savoir s'il en existe un fait
perdre dix minutes en épreuve.
\end{notepourlaclasse}

% =========================================================================
\begin{devoir}[2 heures]
Trois exercices indépendants, notés sur $20$. Le devoir est cumulatif : il porte sur ce chapitre et sur ceux qui le précèdent. Les énoncés sont repris de sessions réelles et assemblés ici en épreuve ; leur provenance est indiquée à chaque fois.

\exods{1}{Baccalauréat, Congo-Brazzaville, série C, session 2022 — exercice 2}{5 points}
$ABC$ est un triangle équilatéral direct de centre de gravité $G$, $O$ est le
milieu de $[AB]$, $I$ le symétrique de $C$ par rapport à $O$, et $BC = 5$~cm.
Soit $(\mathcal C)$ le cercle de diamètre $[IG]$. Montrer que
$(\mathcal C)$ est l'ensemble des points $M$ du plan tels que
$MC = 2MO$.

\exods{2}{Baccalauréat, Madagascar, série C, session 2012 — problème 1, partie I}{7 points}
$ABC$ est un triangle rectangle et isocèle en $A$ tel que $AB = AC = 2$. On
note $G$ le milieu de $[BC]$. Déterminer et construire le barycentre $F$ des
points $A$, $B$ et $C$ affectés respectivement des coefficients $2$, $-1$ et
$1$, puis montrer que $ABGF$ est un parallélogramme.

\exods{3}{Baccalauréat, Mauritanie, série C, session 2022 — exercice 3}{8 points}
$ABC$ est un triangle équilatéral direct de centre $O$, $D$ est le symétrique
de $O$ par rapport à $(AC)$, et $I$, $J$, $K$ sont les milieux respectifs de
$[AB]$, $[BC]$, $[CA]$.
\begin{enumerate}[label=\textbf{\alph*)}, leftmargin=*, itemsep=0.15em]
  \item Montrer qu'il existe un unique déplacement $f$ transformant $B$ en $C$
        et $O$ en $D$.
  \item Justifier que $f$ est une rotation et préciser ses éléments.
  \item Déterminer la nature et les éléments de
        $\trans{AB} \circ f$.
\end{enumerate}

\end{devoir}

\begin{corrige}[devoir surveillé]
\textbf{Exercice 1.} $MC = 2MO$ équivaut à $MC^{2}-4MO^{2} = 0$. La somme des
coefficients vaut $-3$ ; en plaçant $O$ à l'origine et $C$ sur un axe, les
points de la droite $(OC)$ vérifiant la relation sont d'abscisses $-c$ et
$\frac c3$, c'est-à-dire exactement $I$ et $G$. L'ensemble est donc le cercle
de diamètre $[IG]$.

\textbf{Exercice 2.} La somme vaut $2$, donc
$\vect{AF} = \frac12\left(-\vect{AB}+\vect{AC}\right) = \frac12\vect{BC}$. Or
$\vect{BG} = \frac12\vect{BC}$ puisque $G$ est le milieu de $[BC]$. Les
vecteurs $\vect{AF}$ et $\vect{BG}$ sont égaux : $ABGF$ est un
parallélogramme.

\textbf{Exercice 3. a)} $BC = CD$ car $D$ est le symétrique de $O$ par rapport à
$(AC)$ dans une figure équilatérale : les distances sont conservées, donc un
déplacement envoyant $B$ sur $C$ et $O$ sur $D$ existe, et il est unique
puisqu'un déplacement est déterminé par les images de deux points distincts.
\textbf{b)} Ce déplacement n'est pas une translation, car $\vect{BC}$ et
$\vect{OD}$ ne sont pas égaux : c'est une rotation. Son angle est
$\left(\vect{BO},\vect{CD}\right)$ et son centre le point d'intersection des
médiatrices de $[BC]$ et de $[OD]$.
\textbf{c)} $\trans{AB} \circ f$ est un déplacement de même angle que $f$ :
c'est une rotation de même angle, de centre déplacé.

\end{corrige}
